
% Paper F: CC-Image Edge Count Array
\documentclass[11pt,a4paper]{article}
\usepackage[utf8]{inputenc}
\usepackage{amsmath,amssymb,amsthm}
\usepackage{geometry}
\usepackage{booktabs}
\usepackage{array}
\usepackage{hyperref}

\geometry{margin=2.5cm}

\theoremstyle{definition}
\newtheorem{definition}{Definition}[section]
\newtheorem{theorem}{Theorem}[section]
\newtheorem{proposition}{Proposition}[section]
\newtheorem{lemma}{Lemma}[section]
\newtheorem{corollary}{Corollary}[section]
\newtheorem{remark}{Remark}[section]

\title{OEIS Sequence Proposal: CC-Image Edge Count Array\\
\large Measuring the Edge Inflation of the Clique-Cluster Transformation}
\author{Alejandro Zarzuelo Urdiales\\
Department of Mathematics, University of Oviedo\\
UO275823@uniovi.es}
\date{May 2026}

\begin{document}
\maketitle

\begin{abstract}
This paper proposes a new two-parameter integer sequence for the OEIS: the array $T(n, p)$ counting the number of edges in the CC-image of the complete $p$-multigraph on $n$ vertices under the Clique-Cluster Transformation. The array has a piecewise structure with a phase transition at $n = p$, where the diagonal values are perfect cubes $n^3$. We derive closed-form formulas for both regimes, compute values for $n, p \leq 10$, and show that while individual rows and columns overlap with existing OEIS entries (as unrelated polynomial sequences), the unified piecewise array with its graph-theoretic interpretation is entirely absent from the OEIS.
\end{abstract}

\section{Introduction}

The Clique-Cluster Transformation $f_{CC}$ \cite{zarzuelo2025} encodes $p$-multigraphs into simple graphs by replacing each vertex with a coding clique of size $K$ and wiring inter-cluster edges to encode multiplicity. While the vertex count of the CC-image is $n \cdot K$, the edge count determines the true space complexity of the encoding. This paper proposes the two-parameter array of CC-image edge counts for complete $p$-multigraphs and demonstrates its novel combinatorial structure.

\section{Mathematical Framework}

\subsection{CC-Image Edge Count}

\begin{definition}[CC-Image of a Complete $p$-Multigraph]
For the complete $p$-multigraph $K_n^{(p)}$ on $n$ vertices, the CC-transformation with minimum valid coding clique size $K = \max(p, n) + 1$ produces a simple graph $f_{CC}(K_n^{(p)})$ on $n \cdot K$ vertices with:
\begin{enumerate}
\item \textbf{Intra-cluster edges}: Each of the $n$ clusters induces a complete graph $K_K$, contributing $n \cdot \binom{K}{2} = \frac{nK(K-1)}{2}$ edges.
\item \textbf{Inter-cluster edges}: Each of the $\binom{n}{2}$ pairs of clusters is connected by $p$ index-rigid edges (corresponding to multiplicity $p$), contributing $\binom{n}{2} \cdot p = \frac{n(n-1)p}{2}$ edges.
\end{enumerate}
\end{definition}

\begin{definition}[Sequence F: CC-Image Edge Count Array]
For $n \geq 1$ and $p \geq 1$, the array $T(n, p)$ is defined as the total number of edges in the CC-image of $K_n^{(p)}$:
\[
T(n, p) = \frac{nK(K-1)}{2} + \frac{n(n-1)p}{2},
\]
where $K = \max(p, n) + 1$.
\end{definition}

\section{Main Results}

\subsection{Piecewise Closed Form}

\begin{theorem}[Piecewise Formula for $T(n, p)$]
\label{thm:piecewise}
\[
T(n, p) = \begin{cases}
\dfrac{np(p + n)}{2} & \text{if } n \leq p, \\[8pt]
\dfrac{n\bigl[n^2 + (p+1)n - p\bigr]}{2} & \text{if } n > p.
\end{cases}
\]
\end{theorem}

\begin{proof}
\textbf{Case $n \leq p$:} Then $K = p + 1$. Substituting:
\[
T(n, p) = \frac{n(p+1)p}{2} + \frac{n(n-1)p}{2} = \frac{np\bigl[(p+1) + (n-1)\bigr]}{2} = \frac{np(p+n)}{2}.
\]

\textbf{Case $n > p$:} Then $K = n + 1$. Substituting:
\[
T(n, p) = \frac{n(n+1)n}{2} + \frac{n(n-1)p}{2} = \frac{n\bigl[n(n+1) + (n-1)p\bigr]}{2} = \frac{n\bigl[n^2 + n + np - p\bigr]}{2} = \frac{n\bigl[n^2 + (p+1)n - p\bigr]}{2}.
\]
\end{proof}

\subsection{Diagonal Values}

\begin{corollary}[Cubes on the Diagonal]
At $n = p$, both formulas agree and give:
\[
T(n, n) = \frac{n \cdot n \cdot (n + n)}{2} = \frac{2n^3}{2} = n^3.
\]
The diagonal of the array consists of perfect cubes: $1, 8, 27, 64, 125, \ldots$ (OEIS A000578).
\end{corollary}

\begin{remark}[Significance of the Diagonal]
The cube $T(n, n) = n^3$ reflects the balanced case where the maximum multiplicity $p = n$ equals the clique number $\omega(K_n^{(n)}) = n$. At this balance point, the intra-cluster edges contribute $\frac{n^2(n+1)}{2}$ and the inter-cluster edges contribute $\frac{n^2(n-1)}{2}$, summing to $n^3$. The coding clique size is $K = n + 1$ from both directions of the max function.
\end{remark}

\subsection{Structural Properties}

\begin{proposition}[Row Structure for Fixed $n > p$]
For fixed $n$ and $p < n$, the row values form an arithmetic progression with common difference $\frac{n(n-1)}{2}$:
\[
T(n, p) = T(n, 0) + \frac{n(n-1)}{2} \cdot p,
\]
where $T(n, 0) = \frac{n^2(n+1)}{2}$ is the intra-cluster edge count (CC-image of the empty multigraph, formally $p = 0$).
\end{proposition}

\begin{proof}
For $n > p$, the formula is $T(n, p) = \frac{n^2(n+1)}{2} + \frac{n(n-1)p}{2}$. This is linear in $p$ with coefficient $\frac{n(n-1)}{2}$.
\end{proof}

\begin{proposition}[Column Structure for Fixed $p < n$]
For fixed $p$ and large $n$ (specifically $n > p$), the column values follow a cubic polynomial in $n$:
\[
T(n, p) = \frac{n^3}{2} + \frac{(p+1)n^2}{2} - \frac{pn}{2}.
\]
\end{proposition}

\begin{proposition}[Row Structure for Fixed $n \leq p$]
For fixed $n$ and $p \geq n$, the row values are quadratic in $p$:
\[
T(n, p) = \frac{n \cdot p^2}{2} + \frac{n^2 \cdot p}{2}.
\]
\end{proposition}

\begin{corollary}[Asymptotic Edge Inflation Ratio]
For large $n$ with fixed $p$, the CC-image has $\sim n^3/2$ edges while the original complete $p$-multigraph has $\binom{n}{2} \cdot p \sim pn^2/2$ edges. The edge inflation ratio is:
\[
\frac{T(n, p)}{e(K_n^{(p)})} \sim \frac{n^3/2}{pn^2/2} = \frac{n}{p} \to \infty.
\]
For the balanced case $p = n$, both the CC-image and the original have order $n^3$ edges, giving a constant inflation ratio approaching 1 in the leading term.
\end{corollary}

\section{Computational Results}

\subsection{The Array $T(n, p)$ for $n, p = 1, \ldots, 8$}

\begin{center}
\begin{tabular}{c|cccccccc}
\toprule
$n \backslash p$ & 1 & 2 & 3 & 4 & 5 & 6 & 7 & 8 \\
\midrule
1 & 1 & 3 & 6 & 10 & 15 & 21 & 28 & 36 \\
2 & 7 & 8 & 15 & 24 & 35 & 48 & 63 & 80 \\
3 & 21 & 24 & 27 & 42 & 60 & 81 & 105 & 132 \\
4 & 46 & 52 & 58 & 64 & 90 & 120 & 154 & 192 \\
5 & 85 & 95 & 105 & 115 & 125 & 165 & 210 & 260 \\
6 & 141 & 156 & 171 & 186 & 201 & 216 & 273 & 336 \\
7 & 217 & 238 & 259 & 280 & 301 & 322 & 343 & 414 \\
8 & 316 & 344 & 372 & 400 & 428 & 456 & 484 & 512 \\
\bottomrule
\end{tabular}
\end{center}

\subsection{Extended Array for $n, p = 9, 10$}

\begin{center}
\begin{tabular}{c|cc}
\toprule
$n \backslash p$ & 9 & 10 \\
\midrule
1 & 45 & 55 \\
2 & 99 & 120 \\
3 & 162 & 195 \\
4 & 234 & 280 \\
5 & 315 & 380 \\
6 & 405 & 490 \\
7 & 504 & 610 \\
8 & 612 & 740 \\
9 & 729 & 885 \\
10 & 855 & 1000 \\
\bottomrule
\end{tabular}
\end{center}

\subsection{Antidiagonal Reading}

The array read by antidiagonals ($n + p = \text{const}$) gives the sequence:
\[
1, 3, 7, 6, 8, 21, 10, 15, 24, 46, 15, 24, 27, 52, 85, 21, 35, 42, 58, 95, 141, \ldots
\]

This flattened sequence captures the full two-parameter structure in a single linear sequence and is the natural format for OEIS submission of a 2D array.

\section{Why This Is Not a Trivial Extension}

While individual rows and columns of $T(n, p)$ overlap numerically with existing OEIS entries, these overlaps are coincidental --- the existing entries have completely different combinatorial interpretations, and none captures the piecewise structure or the CC-transformation meaning.

\begin{center}
\begin{tabular}{llll}
\toprule
Component & OEIS & Existing meaning & CC meaning? \\
\midrule
Row $n = 1$ & A000217 & Triangular numbers & Single coding clique \\
Column $p = 1$ & A127736 & Composition counts & Simple graph encoding \\
Row $n = 2$ tail & A005563 & $(n+1)^2 - 1$ & Two-cluster encoding \\
Row $n = 3$ tail & A140091 & Hexagonal triangle edges & Three-cluster encoding \\
Diagonal $n = p$ & A000578 & Perfect cubes & Balanced regime \\
\bottomrule
\end{tabular}
\end{center}

The key arguments for non-triviality are:

\begin{enumerate}
\item \textbf{No existing entry captures the piecewise structure.} The transition at $n = p$ from the quadratic regime $np(p+n)/2$ to the cubic regime $n[n^2 + (p+1)n - p]/2$ is a fundamental feature of the CC-transformation that no single polynomial sequence can represent.

\item \textbf{No existing entry has the CC-transformation interpretation.} The graph-theoretic meaning --- edge inflation of a multigraph-to-simple-graph encoding --- is entirely absent from the OEIS. None of the overlapping entries (A127736, A140091) mention multigraphs, coding cliques, or graph embeddings.

\item \textbf{The 2D array as a unified object is novel.} While each row for fixed $n$ is eventually a polynomial in $p$ (and vice versa), the \emph{array as a whole} with its phase transition and graph-theoretic meaning does not exist in the OEIS. The antidiagonal reading $1, 3, 7, 6, 8, 21, 10, 15, 24, 46, \ldots$ has no OEIS match.

\item \textbf{Individual columns for $p \geq 3$ are novel.} The sequences for fixed $p \geq 3$ with the piecewise first terms are not in the OEIS. For $p = 3$: the values $6, 15, 27, 58, 105, 171, 259, 372, 513, 685, \ldots$ have no OEIS match.
\end{enumerate}

\section{OEIS Cross-References}

\begin{itemize}
\item \textbf{A000217}: Triangular numbers $p(p+1)/2$. Row $n = 1$ of our array.
\item \textbf{A127736}: $a(n) = n(n^2 + 2n - 1)/2$. Column $p = 1$ of our array (without graph interpretation).
\item \textbf{A005563}: $a(n) = n(n+2)$. Row $n = 2$ for $p \geq 2$ (without graph interpretation).
\item \textbf{A140091}: $a(n) = 3n(n+3)/2$. Row $n = 3$ for $p \geq 3$ (hexagonal triangle edges, without graph interpretation).
\item \textbf{A000578}: Perfect cubes $n^3$. Diagonal of our array.
\item \textbf{A002378}: Oblong numbers $n(n+1)$. Vertex count of CC-images (related but different parameter).
\end{itemize}

\section{Algorithmic Implementation}

\begin{verbatim}
def T(n, p):
    """Compute the CC-Image edge count T(n, p)."""
    K = max(p, n) + 1
    intra = n * K * (K - 1) // 2
    inter = n * (n - 1) * p // 2
    return intra + inter

# Or using the piecewise formula:
def T_piecewise(n, p):
    if n <= p:
        return n * p * (p + n) // 2
    else:
        return n * (n**2 + (p + 1) * n - p) // 2

# Compute the array
for n in range(1, 11):
    row = [T(n, p) for p in range(1, 11)]
    print(f"n={n}: {row}")

# Compute antidiagonal reading
antidiag = []
for s in range(2, 22):
    for n in range(1, s):
        p = s - n
        if p >= 1:
            antidiag.append(T(n, p))
print(f"Antidiagonal: {antidiag}")
\end{verbatim}

\section{Applications}

\begin{enumerate}
\item \textbf{Space Complexity of Graph Encoding}: The array $T(n, p)$ directly measures the edge-space cost of encoding a $p$-multigraph as a simple graph via the CC-transformation. The ratio $T(n, p) / e(G)$ quantifies the encoding overhead.

\item \textbf{Algorithmic Feasibility}: The cubic growth $T(n, p) \sim n^3/2$ for fixed $p$ shows that the CC-transformation is polynomial in the input size, making it feasible for computational applications. The balanced case $p = n$ gives $T = n^3$, which matches the $\Theta(n^3)$ complexity of matrix multiplication --- suggesting a potential connection.

\item \textbf{Graph Database Design}: In applications where multigraphs must be stored in simple-graph databases (e.g., knowledge graphs in triple stores), the array $T(n, p)$ gives the exact storage overhead as a function of the original graph's density and multiplicity.
\end{enumerate}

\section{Conclusion}

The CC-Image Edge Count Array $T(n, p)$ provides the first complete characterization of the edge inflation produced by the Clique-Cluster Transformation. Its piecewise structure with a phase transition at $n = p$, the perfect-cube diagonal, and the graph-theoretic interpretation distinguish it from all existing OEIS entries. The antidiagonal reading offers a natural linear format for OEIS inclusion.

\begin{thebibliography}{9}
\bibitem{zarzuelo2025} A. Zarzuelo Urdiales, ``Canonical Embedding of Universal $p$-Graphs, Commensurable Weighted Graphs, and Line Structures into the Rado Graph via Clique-Size Dominance Transformation,'' December 2025.
\bibitem{oeis} N. J. A. Sloane et al., The On-Line Encyclopedia of Integer Sequences, \url{https://oeis.org}.
\bibitem{diestel2017} R. Diestel, \textit{Graph Theory}, 5th ed., Springer, 2017.
\end{thebibliography}

\end{document}
